\section{Problemfelt}

Viborg Katedralskole fylder 100 år i 2026. I forbindelse med jubilæet ønsker skolen at udvikle en digital oplevelse, der kan inddrage nuværende elever og formidle bygningens historie på.

Målgruppen er elever i alderen 14-19 år. Forskning viser, at denne aldersgruppe engageres, når indhold opleves som relevant for deres egen livsverden, og når de får mulighed for at bidrage aktivt \cite{VanDoorsselaere2025}. Digitale oplevelser kan øge motivationen, men effekten afhænger af designet og den didaktiske ramme \cite{Wouters2013, Mayer2022}.

Skolen har i dag fysiske mindesmærker og årlige ritualer,\cite{katedralskole2026} men ingen interaktiv digital platform, der gør historien nærværende for de unge. Der mangler viden om, hvordan en bygning med 100 års historie kan fungere som fundament for en digital oplevelse, der både respekterer autenticiteten og skaber ejerskab hos eleverne. \\

Projektet undersøger derfor, hvordan Viborg Katedralskoles historie kan anvendes som fundament for udviklingen af en digital oplevelse i forbindelse med skolens jubilæum, samt hvordan nuværende elever kan inddrages i processen for at skabe engagement og interesse omkring skolens historiske fortællinger.


% Efter den 9. april 1940, var Danmark besat af Tyskland under 2. Verdenskrig \cite{BRUNBECH2025}, dette betød for mange dansker at hverdagen var hård og trist, grundet ” varemangel, rationering, luftalarmer, regler om mørklægning om aftenen og lukkede grænser"\cite{nationalmuseet2026}, efter cirka et års besættelse, fik flere Danskere nok af undertrykkelsen, og der begyndte at dukke Danske modstandsbevægelser op flere steder i Danmark, disse modstands bevægelser fik blandt andet hjælp fra det Britiske Special Operations Executive (SOE) som hjalp flere modstands folk i Danmark med forskelige former for oprør og propaganda\cite{trommer2026}. I Viborg dukkede der også flere modstandsbevægelser op (Fabricius Wøldike, 2026)\cite{Woldlike2026}, blandt disse modstandsfolk i Viborg, var der elever fra Viborg katedral Skole, som fik store ansvarsposter under modstandskampen. En af aktiviteterne som bliv udført var at distribuere Viborg Stifts tiden, som på dette tidspunkt var ulovligt og fyldt med propaganda, dette var den tyske ordensmagt anført af Gestapo ikke tilfredse med og det kulminerede i en razzia på Viborg katedral skole, den 23 oktober 1944, hvor 14 modstands folk fra Viborg blev arresteret, disse medlemmer også med dybe spor i modstandsbevægelserne, dette betød blandt andet af Søren Vestergaard måtte flygte og gemme sig fra byen, for ikke at blive fanget, de illegale aviser, fik gestapo dog ikke gjort op med, da 2 elever fra skolen over tog arbejdet med udgivelsen af den ulovlige avis.\cite{katedralskole2026}. Flere elever fra skolen Fik også andre store roller i kampen for frihed, bland andet kom Christian Ulrik Hansen også fra Viborg Katedralskole\cite{Damgaard1990}, ham samt andre danske modstands Folk fra Viborg af, er nævnt på mindetavlen uden for skolen, til minde om de ofre de gav for befri Danmark, i dag mindes de, hvert år på befrielse-dagen\cite{viborgmuseum2026}.
% har brug for meget mere